% Generated semantic source for AI Almanach v3.
% Edit v3 directly after initial generation; do not overwrite
% editorial changes by rerunning the converter casually.

% ==== Foundational Computer Vision =============================================================================
\chapter{Foundational Computer Vision}

\chapterlead{How do pixels and geometry become testable statements about a scene?}{Sense}{Represent}{Infer}{Stress-test}{computer-vision}

% ==== Section 1: Image Processing and Low-Level Vision ==========================================
\label{ch:computer-vision}

\begin{learningobjectives}
After working through this chapter you should be able to:
\begin{itemize}
  \item use the pinhole model to predict how many pixels tall an object will be
        at a given distance --- and therefore whether it is detectable at all;
  \item apply a Sobel kernel by hand and read gradient magnitude and direction
        from the result;
  \item compute intersection-over-union and explain why the same detection
        passes at one threshold and fails at another;
  \item calculate the receptive field of a stack of convolutions and justify
        small kernels from the parameter count;
  \item estimate activation memory for a training batch and explain why it,
        not the weights, sets your batch size;
  \item name the failure mode a vision model exhibits under distribution shift
        and the test that would have caught it.
\end{itemize}
\end{learningobjectives}

\begin{plainlanguage}
A camera destroys information on purpose. It flattens a three-dimensional world
onto a grid of numbers, throwing away depth, and it samples that grid at finite
resolution, throwing away detail. Computer vision is the business of putting
back what the projection removed --- using geometry where geometry suffices,
and learned priors where it does not. Almost every hard problem in this chapter
is a consequence of that first, deliberate loss.
\end{plainlanguage}

\begin{engineernote}
Vision systems fail on the boundary conditions of the sensor far more often
than on the cleverness of the model: too few pixels on target, motion blur,
glare, rain, an exposure setting that clips. Before improving an architecture,
compute how many pixels the object of interest actually occupies at the range
you care about. If the answer is twenty, no model will save you --- but a
different lens will.
\end{engineernote}

\section{Image Processing and Low-Level Vision}

% ---- 9.1.1 Image Acquisition ------------------------------------------------------------------
\subsection{Image Acquisition}

\begin{workedexample}[How many pixels is a pedestrian?]
A camera with focal length $f = 8\,\text{mm}$ and a sensor whose pixels are
$4\,\upmu\text{m}$ across. Under the pinhole model, an object of height $H$ at
distance $Z$ projects to image height
\[
  h = f\,\frac{H}{Z} .
\]

\emph{A $1.7\,\text{m}$ pedestrian at $20\,\text{m}$:}
\[
  h = 8\,\text{mm}\times\frac{1700\,\text{mm}}{20{,}000\,\text{mm}}
    = 0.68\,\text{mm}
  \quad\Rightarrow\quad
  \frac{0.68}{0.004} = 170 \text{ pixels tall.}
\]
Comfortable for any detector.

\emph{The same pedestrian at $100\,\text{m}$:}
\[
  h = 8\times\frac{1700}{100{,}000} = 0.136\,\text{mm}
  \quad\Rightarrow\quad 34 \text{ pixels tall.}
\]

\emph{Why this matters.} Detection performance degrades sharply below roughly
$30$ pixels of object height. At $100\,\text{m}$ this camera is at the edge of
being able to see a person at all --- and a vehicle travelling at
$30\,\text{m/s}$ covers $100\,\text{m}$ in $3.3$ seconds. The sensing range,
not the network, sets the safety margin. Note also that $h \propto 1/Z$: doubling
the detection distance requires doubling the focal length, which halves the
field of view. That trade-off cannot be trained away.
\end{workedexample}

\begin{v3topic}{The Human Visual System}
    The human retina contains approximately \textbf{120 million rods} (scotopic: low-light, achromatic) and
    \textbf{6 million cones} (photopic: three types L/M/S sensitive to long, medium, short wavelengths)\textsuperscript{\cite[][p.~57]{szeliski2022}}.
    The \textbf{fovea} packs high-density cones for fine detail; peripheral vision relies mostly on rods.
    Lateral inhibition in the retina enhances local contrast before signals reach the visual cortex.
\visualseparator
    CV relevance: the trichromatic cone response motivates RGB color models; the multi-scale,
    hierarchical organisation of V1--V5 cortical areas inspired convolutional neural networks.

\end{v3topic}
\begin{v3topic}{Camera Sensors: CCD and CMOS}
    Most digital cameras capture light with one of two sensor technologies\textsuperscript{\cite[][pp.~68--73]{jahneDigitalImageProcessing2005}}:
    \begin{itemize}
        \item \textbf{CCD} (Charge-Coupled Device): charge shifts pixel-by-pixel to a single readout amplifier; excellent uniformity, higher power consumption.
        \item \textbf{CMOS}: each pixel has its own amplifier enabling random access; lower power, faster readout, now dominant.
    \end{itemize}
    Most sensors use a \textbf{Bayer mosaic} (one R, two G, one B per $2{\times}2$ block).
    Colour reconstruction via \textbf{demosaicing}.
    Principal noise sources: shot noise (Poisson), read noise (Gaussian), fixed-pattern noise.
\end{v3topic}

% ---- 9.1.2 Image Representation and Morphology ------------------------------------------------
\subsection{Image Representation and Morphology}

\begin{v3topic}{Image Types and Colour Spaces}
    A digital image is a discrete function $I: \mathbb{Z}^2 \to V$ where $V$ is the value set\textsuperscript{\cite[][p.~38]{gonzalezWoods2018}}.
    Common types: \textbf{binary} (0/1), \textbf{grayscale} ($0$--$255$), \textbf{colour} (3-channel), \textbf{depth/range} (float).
\visualseparator
    \begin{tblr}{width=\linewidth,colspec={|Q[l,m,2cm]|X[l]|X[l]|}}
        \hline
        \textbf{Space} & \textbf{Channels} & \textbf{Key use} \\
        \hline
        RGB   & Red, Green, Blue (additive) & Capture, display; device-dependent \\
        \hline
        HSV   & Hue, Saturation, Value & Illumination-robust hue; segmentation \\
        \hline
        CIE L*a*b* & Lightness, red-green, blue-yellow & Perceptually uniform; colour difference $\Delta E$ \\
        \hline
        YCbCr & Luma + two chrominance channels & JPEG/video compression \\
        \hline
    \end{tblr}
\captionof{table}{Colour spaces and what each is good for. Converting to a space that separates intensity from colour makes many operations robust to lighting changes.}
\label{tab:cv-colour-spaces}

\begin{workedexample}[Why HSV survives a change in lighting and RGB does not]
Take an orange pixel, $\text{RGB} = (200, 100, 50)$. Normalise to $[0,1]$:
$R=0.784$, $G=0.392$, $B=0.196$.

\emph{Convert to HSV.} With $\max = 0.784$, $\min = 0.196$, and
$\Delta = 0.588$:
\[
\begin{aligned}
  V &= \max = 0.784, \\
  S &= \Delta/\max = 0.588/0.784 = 0.750, \\
  H &= 60^\circ \times \frac{G-B}{\Delta}
     = 60^\circ \times \frac{0.196}{0.588} = 20^\circ
     \quad(\text{since } \max = R).
\end{aligned}
\]

\emph{Now move the object into shade}, halving every channel:
$\text{RGB} = (100, 50, 25)$, so $R=0.392$, $G=0.196$, $B=0.098$,
$\max=0.392$, $\Delta=0.294$:
\[
  V = 0.392,\qquad
  S = \frac{0.294}{0.392} = 0.750,\qquad
  H = 60^\circ \times \frac{0.098}{0.294} = 20^\circ .
\]

\emph{Read the result.} \textbf{All three RGB values changed; only $V$
changed in HSV.} Hue and saturation are exactly as they were, because the
illumination change was a common scale factor that $S$ and $H$ both divide
out.

\emph{The engineering consequence.} A colour-based segmentation rule written
in RGB (``red is above 180'') fails as soon as a cloud passes. The same rule
written as ``hue between $10^\circ$ and $30^\circ$'' survives it. That is the
entire reason HSV exists in a vision pipeline --- and also its limit: the
division by $\max$ becomes numerically unstable for very dark pixels, where
hue is meaningless.
\end{workedexample}

\end{v3topic}
\begin{v3topic}{Morphological Operations}
    Morphology operates on binary (or grayscale) images using a \textbf{structuring element} $B$\textsuperscript{\cite[][pp.~627--636]{gonzalezWoods2018}}:
    \begin{itemize}
        \item \textbf{Erosion} $A \ominus B$: pixel belongs to result only if all of $B$ fits inside $A$; shrinks objects, removes small blobs.
        \item \textbf{Dilation} $A \oplus B$: pixel belongs to result if any part of $B$ overlaps $A$; expands objects, fills thin gaps.
    \end{itemize}
    Grayscale morphology substitutes min (erosion) and max (dilation) over the structuring element neighbourhood.

\end{v3topic}
\begin{v3topic}{Opening and Closing}
    Compositions of erosion and dilation yield powerful filters\textsuperscript{\cite[][pp.~638--645]{gonzalezWoods2018}}:
    \begin{itemize}
        \item \textbf{Opening} $= A \ominus B \oplus B$: smooth contour, remove small protrusions and isolated noise pixels.
        \item \textbf{Closing} $= A \oplus B \ominus B$: fill small holes and narrow gaps between objects.
    \end{itemize}
    The \textbf{top-hat transform} $T = I - \mathrm{Opening}(I)$ extracts bright features smaller than the structuring element — useful for illumination correction.
\end{v3topic}

% ---- 9.1.3 Single and Multi-View Geometry -----------------------------------------------------
\subsection{Single and Multi-View Geometry}

\begin{v3topic}{Pinhole Camera Model}
    The \textbf{pinhole camera} models perspective projection\textsuperscript{\cite[][pp.~44--50]{hartleyZisserman2004}}:
    \begin{equation}
        \lambda \begin{bmatrix} u \\ v \\ 1 \end{bmatrix}
        = \underbrace{\begin{bmatrix} f_x & 0 & c_x \\ 0 & f_y & c_y \\ 0 & 0 & 1 \end{bmatrix}}_{K}
        \begin{bmatrix} R \mid \mathbf{t} \end{bmatrix}
        \begin{bmatrix} X \\ Y \\ Z \\ 1 \end{bmatrix}
    \end{equation}
\where{$(X,Y,Z)$ is a point in world coordinates and $(u,v)$ its pixel
position; $K$ is the intrinsic matrix with focal lengths $f_x,f_y$ in pixels and
principal point $(c_x,c_y)$; $R$ and $\mathbf{t}$ are the camera's rotation and
translation, and $\lambda$ the depth scale homogeneous coordinates leave free.}
    $K$ = intrinsic matrix ($f_x, f_y$: focal lengths in pixels; $c_x, c_y$: principal point).
    $[R \mid \mathbf{t}]$: extrinsic (rotation + translation from world to camera).
\visualseparator
    \begin{center}
    \begin{tikzpicture}[scale=0.85,
        arr/.style={-{Stealth[length=2mm]}, thick},
        ray/.style={dashed, gray}
    ]
        % Image plane
        \draw[thick, fill=blue!8] (-0.5, -0.8) rectangle (0.5, 0.8);
        \node[font=\scriptsize] at (0, 1.0) {Image plane};
        % Optical centre
        \filldraw[black] (2.0, 0) circle (2pt) node[above, font=\scriptsize] {$C$};
        % 3D point
        \filldraw[orange] (4.0, 1.2) circle (2pt) node[right, font=\scriptsize] {$P(X,Y,Z)$};
        % Project onto image plane
        \draw[ray] (4.0, 1.2) -- (2.0, 0) -- (0.3, 0.45);
        \filldraw[red] (0.3, 0.45) circle (1.5pt) node[left, font=\scriptsize] {$p(u,v)$};
        % Optical axis
        \draw[arr, thin] (2.0, 0) -- (-0.5, 0) node[midway, below, font=\scriptsize] {$f$};
        % Z axis
        \draw[arr, thin, gray] (2.0,0) -- (4.5, 0) node[right, font=\scriptsize] {$Z$};
    \end{tikzpicture}
\captionof{figure}{The pinhole camera model. Image height scales as $f H / Z$,
so object size in pixels falls off as $1/Z$ --- the relationship that decides
whether a detector can see anything at all at range.}
\label{fig:cv-pinhole}
    \end{center}

\end{v3topic}
\begin{v3topic}{Homography and Epipolar Geometry}
    A \textbf{homography} $H$ (3×3, 8 DOF) maps points between two views of a planar surface\textsuperscript{\cite[][pp.~32--36]{hartleyZisserman2004}}:
    $\mathbf{p}' \sim H\mathbf{p}$.
    Applications: panorama stitching, planar marker tracking.
    Estimated robustly with \textbf{RANSAC} + DLT from $\geq 4$ point correspondences.
\visualseparator
    \textbf{Epipolar geometry} describes the constraint between two calibrated views.
    For a point $\mathbf{p}$ in image 1, its match in image 2 lies on the
    \textbf{epipolar line} $\ell' = F\mathbf{p}$, where $F$ is the $3{\times}3$ \textbf{fundamental matrix} (rank 2).
    Stereo depth: $Z = \frac{f \cdot B}{d}$ (baseline $B$, disparity $d = x_L - x_R$).
\end{v3topic}

% ---- 9.1.4 Filtering --------------------------------------------------------------------------
\subsection{Filtering}

\begin{v3topic}{Spatial Filtering and Convolution}
    \textbf{Convolution} with kernel $h$ over image $f$ at position $(x,y)$\textsuperscript{\cite[][pp.~120--132]{gonzalezWoods2018}}:
    \begin{equation}
        (f * h)[x,y] = \sum_{i,j} f[x-i,\, y-j]\, h[i,j]
    \end{equation}
\where{$f$ is the image, $h$ the kernel, $(x,y)$ the output pixel, and $(i,j)$
range over the kernel. Note the minus signs: true convolution flips the kernel,
whereas most deep-learning libraries implement cross-correlation without the
flip.}
    Common spatial filters:
    \begin{itemize}
        \item \textbf{Box filter}: uniform kernel; fast mean blur, but can create blocky smoothing and boundary artefacts.
        \item \textbf{Gaussian}: smooth and isotropic; $\sigma$ controls blur radius and typically avoids the ringing associated with sharp frequency cut-offs.
        \item \textbf{Unsharp mask}: $I_{\text{sharp}} = I + \alpha(I - G_\sigma * I)$; amplifies edges.
        \item \textbf{Median filter}: non-linear; preserves edges; best for salt-and-pepper noise.
    \end{itemize}

\end{v3topic}
\begin{v3topic}{Gaussian Blur}
    The \textbf{Gaussian kernel} is the most important smoothing filter\textsuperscript{\cite[][pp.~138--142]{jahneDigitalImageProcessing2005}}:
    \begin{equation}
        G(x,y;\sigma) = \frac{1}{2\pi\sigma^2} \exp\!\left(-\frac{x^2+y^2}{2\sigma^2}\right)
    \end{equation}
\where{$\sigma$ sets the blur radius --- larger $\sigma$ removes more detail and
more noise --- and the leading factor normalises the kernel so overall brightness
is unchanged.}
    Key properties:
    \begin{itemize}
        \item Separable: $G(x,y) = G(x)\cdot G(y)$ — fast 1D × 1D implementation.
        \item Gaussian scale space: progressively blurring at increasing $\sigma$ reveals coarser structure; foundation of multi-scale analysis and SIFT.
        \item Isotropic: no preferred orientation → rotationally invariant smoothing.
    \end{itemize}

\end{v3topic}
\begin{v3topic}{Frequency Domain Filtering}
    The \textbf{Discrete Fourier Transform} decomposes an image into spatial frequency components\textsuperscript{\cite[][pp.~218--245]{gonzalezWoods2018}}:
    \begin{equation}
        F(u,v) = \sum_{x=0}^{M-1}\sum_{y=0}^{N-1} f(x,y)\, e^{-j2\pi\left(\frac{ux}{M}+\frac{vy}{N}\right)}
    \end{equation}
\where{$f(x,y)$ is the image of size $M\times N$, $(u,v)$ the horizontal and
vertical spatial frequencies, and $j$ the imaginary unit; $|F(u,v)|$ is the
magnitude spectrum usually displayed.}
    The \textbf{convolution theorem} states: $(f * h) \leftrightarrow F \cdot H$ — convolution in the spatial domain equals pointwise multiplication in the frequency domain.
    This allows large kernel convolutions to be performed efficiently via FFT in $O(MN\log MN)$.
\visualseparator
    \begin{tblr}{width=\linewidth,colspec={|Q[l,m]|X[l]|X[l]|}}
        \hline
        \textbf{Filter} & \textbf{What it passes} & \textbf{Effect} \\
        \hline
        Low-pass  & Low frequencies & Blurring, noise removal \\
        \hline
        High-pass & High frequencies & Edge enhancement, sharpening \\
        \hline
        Band-pass & Specific range & Texture analysis, Gabor filters \\
        \hline
    \end{tblr}
\captionof{table}{Frequency-domain filters. Low-pass smooths and blurs; high-pass sharpens and amplifies noise --- the same trade in both directions.}
\label{tab:cv-frequency-filters}
\end{v3topic}

% ---- 9.1.5 Texture ----------------------------------------------------------------------------
\subsection{Texture}

\begin{v3topic}{Classical Texture Descriptors}
    Hand-crafted texture features capture local intensity patterns\textsuperscript{\cite[][pp.~645--662]{szeliski2022}}:
    \begin{itemize}
        \item \textbf{GLCM} (Gray-Level Co-occurrence Matrix): counts how often pairs of pixel intensities co-occur at a given offset; features: contrast, energy, homogeneity, correlation.
        \item \textbf{LBP} (Local Binary Pattern): compare centre pixel to its 8 neighbours; encode as binary string $\to$ histogram. Fast, robust to monotonic illumination changes.
        \item \textbf{Gabor filters}: oriented sinusoidal wave modulated by a Gaussian envelope; tunable frequency $\omega_0$ and orientation $\theta$. A bank at multiple scales and orientations captures rich texture information.
    \end{itemize}

\end{v3topic}
\begin{v3topic}{Bag-of-Visual-Words and CNN Texture}
    \textbf{Bag-of-Visual-Words (BoVW)}\textsuperscript{\cite[][pp.~682--686]{szeliski2022}}: (1) extract local descriptors (SIFT) densely; (2) cluster into a visual codebook ($K$-means, $K$ = 500–4000); (3) represent each image as a frequency histogram over codewords; (4) classify with SVM (TF-IDF weighting).
\visualseparator
    \textbf{CNN texture features}: early layers learn Gabor-like oriented filters; intermediate layers capture textons and repetitive patterns.
    The \textbf{Gram matrix} $G = F^\top F$ (where $F$ are spatial feature maps flattened) captures second-order statistics of CNN activations --- used for neural style transfer and texture synthesis\textsuperscript{\cite[][p.~119]{fleuretLittleBookDeep}}.
\end{v3topic}


% ==== Section 2: Mid-Level Vision and Video ====================================================
\section{Mid-Level Vision and Video}

% ---- 9.2.1 Edge and Feature Detection ---------------------------------------------------------
\subsection{Edge and Feature Detection}

\begin{v3topic}{Edge Detection and the Canny Detector}
    Edges mark discontinuities in intensity caused by depth, surface orientation, or reflectance changes.
    The \textbf{Sobel operator} approximates image gradients\textsuperscript{\cite[][pp.~693--706]{gonzalezWoods2018}}:
    \begin{equation}
        |\nabla I| = \sqrt{G_x^2 + G_y^2}, \quad
        G_x = \begin{bmatrix}-1&0&1\\-2&0&2\\-1&0&1\end{bmatrix} * I
        \label{eq:cv-sobel}
    \end{equation}
\where{$I$ is the image, $G_x$ and $G_y$ the horizontal and vertical gradient
responses from the kernel shown and its transpose, and $|\nabla I|$ the edge
strength; $\operatorname{atan2}(G_y,G_x)$ gives the edge orientation.}

\begin{workedexample}[Sobel on a vertical edge, by hand]
Take a $3\times3$ patch straddling a vertical intensity step:
\[
  I = \begin{pmatrix} 10 & 10 & 50\\ 10 & 10 & 50\\ 10 & 10 & 50 \end{pmatrix}.
\]

\emph{Horizontal gradient} with the $G_x$ kernel above, multiplying
element-wise and summing row by row:
\[
  G_x = \underbrace{(-10+50)}_{\text{row 1}} + \underbrace{(-20+100)}_{\text{row 2}}
      + \underbrace{(-10+50)}_{\text{row 3}} = 40+80+40 = 160 .
\]

\emph{Vertical gradient} with $G_y$ (the same kernel transposed):
\[
  G_y = (-10-20-50) + 0 + (10+20+50) = -80 + 80 = 0 .
\]

\emph{Magnitude and direction.}
\[
  |\nabla I| = \sqrt{160^2 + 0^2} = 160,
  \qquad
  \theta = \operatorname{atan2}(0,160) = 0^\circ .
\]

\emph{Read it.} A strong response, pointing purely in $x$. The gradient
direction is perpendicular to the edge, so a gradient at $0^\circ$ means a
\emph{vertical} edge --- a distinction worth getting right once and never
having to re-derive. And note $G_y = 0$ exactly: the operator is blind to the
edge along its own direction, which is why both components are always needed.
\end{workedexample}
    The \textbf{Canny detector} (1986) is the standard pipeline:
    \begin{enumerate}
        \item Gaussian smoothing ($\sigma$)
        \item Gradient magnitude \& direction
        \item Non-maximum suppression (thin edges to 1 pixel)
        \item Double threshold (high $T_h$, low $T_l$)
        \item Hysteresis linking: strong edges kept; weak edges kept only if connected to strong ones
    \end{enumerate}

\end{v3topic}
\begin{v3topic}{Harris Corners and FAST}
    \textbf{Harris corner detector} (Harris \& Stephens, 1988)\textsuperscript{\cite[][pp.~215--220]{szeliski2022}}:
    build the \textbf{structure tensor} from image gradients:
    \begin{equation}
        M = \sum_{(x,y) \in W} w(x,y) \begin{bmatrix} I_x^2 & I_xI_y \\ I_xI_y & I_y^2 \end{bmatrix}
    \end{equation}
\where{$W$ is a window around the candidate corner, $w(x,y)$ a weighting
(often Gaussian), and $I_x,I_y$ the image gradients. The eigenvalues of $M$
decide the verdict: both large means a corner, one large means an edge.}
    Corner response: $R = \det(M) - k\,\mathrm{trace}^2(M)$.
    Large $R$ $\to$ corner; $R < 0$ $\to$ edge; $|R|$ small $\to$ flat region.
\visualseparator
    \textbf{FAST} (Features from Accelerated Segment Test): test 16 pixels on a circle of radius 3; a pixel is a corner if $\geq N$ consecutive pixels are all brighter or all darker than the centre by a threshold.
    Extremely fast; widely used in real-time SLAM and tracking pipelines.

\end{v3topic}
\begin{v3topic}{SIFT, SURF, and ORB}
    Scale- and rotation-invariant local feature descriptors\textsuperscript{\cite[][pp.~225--240]{szeliski2022}}:
\visualseparator
    \begin{tblr}{width=\linewidth,colspec={|Q[l,m,1.5cm]|X[l]|X[l]|X[l]|}}
        \hline
        \textbf{Method} & \textbf{Detection} & \textbf{Descriptor} & \textbf{Properties} \\
        \hline
        SIFT  & DoG (Difference of Gaussians) multi-scale & 128-dim gradient histogram & Scale + rotation invariant; patented until 2020 \\
        \hline
        SURF  & Hessian approximation with integral images & 64-dim Haar wavelet response & Faster than SIFT; approximate \\
        \hline
        ORB   & FAST keypoints + Harris score & 256-bit BRIEF binary descriptor & Real-time; patent-free; good for mobile/SLAM \\
        \hline
    \end{tblr}
\captionof{table}{Classical keypoint detectors and descriptors, with the invariances each provides. ORB is the one without patent constraints and is fast enough for real time.}
\label{tab:cv-keypoints}
    \textbf{Feature matching}: nearest-neighbour in descriptor space; \textbf{Lowe's ratio test} ($d_1/d_2 < 0.8$) removes ambiguous matches. \textbf{RANSAC} robustly estimates a geometric model (homography) from potentially noisy correspondences.
\end{v3topic}

% ---- 9.2.2 Segmentation -----------------------------------------------------------------------
\subsection{Segmentation}

\begin{v3topic}{Thresholding and Region-Based Segmentation}
    \textbf{Segmentation} partitions an image into meaningful regions\textsuperscript{\cite[][pp.~741--804]{gonzalezWoods2018}}.
    \begin{itemize}
        \item \textbf{Otsu thresholding}: analytically finds global threshold $T^*$ that maximises inter-class variance between foreground and background; optimal for bimodal histograms.
        \item \textbf{Adaptive thresholding}: local threshold based on neighbourhood mean/Gaussian; handles uneven illumination.
        \item \textbf{Watershed}: treat gradient magnitude as topographic surface; region growing from seed markers; boundaries form at watershed lines. Over-segmentation controlled by marker selection.
    \end{itemize}

\end{v3topic}
\begin{v3topic}{Graph-Cut and Superpixels}
    \textbf{Graph-Cut segmentation}\textsuperscript{\cite[][pp.~302--306]{szeliski2022}}: model image as graph ($G = \langle V, E \rangle$) with nodes as pixels and edge weights encoding similarity; minimise the cut to obtain two-class partition.
    \textbf{GrabCut} (Rother et al., 2004): iterative refinement using GMMs for foreground/background appearance, solved via min-cut; user supplies rough bounding box.
\visualseparator
    \textbf{Superpixels} (SLIC, 2012): group pixels into perceptually uniform, compact clusters via $k$-means in $[\mathrm{CIELAB},\,x,\,y]$ with a spatial compactness term.
    Reduce pixels by ${\times}100$ while preserving boundaries; used as pre-processing for semantic segmentation.
\end{v3topic}

% ---- 9.2.3 Motion and Optical Flow ------------------------------------------------------------
\subsection{Motion and Optical Flow}

\begin{v3topic}{The Optical Flow Equation}
    \textbf{Optical flow} $\mathbf{u}=(u,v)$ describes apparent pixel motion between consecutive frames.
    The \textbf{brightness constancy assumption} $I(x,y,t) = I(x+u, y+v, t+\Delta t)$ leads via Taylor expansion to the \textbf{optical flow constraint equation}\textsuperscript{\cite[][pp.~498--504]{szeliski2022}}:
    \begin{equation}
        I_x\,u + I_y\,v + I_t = 0
    \end{equation}
    where $I_x, I_y$ are spatial and $I_t$ the temporal gradient.
    This is one equation in two unknowns --- the \textbf{aperture problem}: flow along the gradient direction is determined; motion parallel to an edge is ambiguous.
    Additional constraints (local smoothness or spatial neighbourhood) resolve the ambiguity.

\end{v3topic}
\begin{v3topic}{Lucas-Kanade and Horn-Schunck}
    Two classical approaches resolve the aperture problem with different assumptions\textsuperscript{\cite[][pp.~504--510]{szeliski2022}}:
    \begin{itemize}
        \item \textbf{Lucas-Kanade (1981)}: local method; assumes constant flow within a $w{\times}w$ window; over-determined system solved by least squares.
          Works well for small displacements; fails in textureless regions.
        \item \textbf{Horn-Schunck (1981)}: global method; adds quadratic smoothness regulariser $\lambda\iint(|\nabla u|^2 + |\nabla v|^2)\,\mathrm{d}x\,\mathrm{d}y$ to the data term; solved iteratively.
          Dense flow; tends to smooth discontinuities at motion boundaries.
    \end{itemize}

\end{v3topic}
\begin{v3topic}{Deep Optical Flow}
    Deep learning overcomes limitations of classical methods on large displacements\textsuperscript{\cite[][p.~83]{fleuretLittleBookDeep}}:
\visualseparator
    \begin{tblr}{width=\linewidth,colspec={|Q[l,m,2cm]|X[l]|X[l]|}}
        \hline
        \textbf{Method} & \textbf{Key Idea} & \textbf{Contribution} \\
        \hline
        FlowNet (2015)  & End-to-end CNN on synthetic ``Flying Chairs'' dataset & First CNN for optical flow; correlation layer \\
        \hline
        PWC-Net (2018)  & Pyramid + warping + cost volume & Efficient; strong accuracy–speed tradeoff \\
        \hline
        RAFT (2020)     & 4D all-pairs cost volume + iterative GRU updates & State-of-the-art; robust to large motions \\
        \hline
    \end{tblr}
\captionof{table}{Learned optical-flow methods and the idea each contributed.}
\label{tab:cv-optical-flow}
    Deep methods generalise better to large motions and handle textureless regions via learned priors; they require large synthetic training datasets (FlyingChairs, Sintel, KITTI).
\end{v3topic}

% ---- 9.2.4 Tracking ---------------------------------------------------------------------------
\subsection{Tracking}

\begin{v3topic}{Kalman Filter: Predict-Update Cycle}
    The \textbf{Kalman filter} (1960) is the optimal linear estimator for tracking under Gaussian noise\textsuperscript{\cite[][pp.~540--547]{szeliski2022}}:
    \begin{itemize}
        \item \textbf{Predict}: $\hat{\mathbf{x}}_{k|k-1} = F\hat{\mathbf{x}}_{k-1}$; $P_{k|k-1} = FPF^\top + Q$
        \item \textbf{Update}: $K = P H^\top (HPH^\top + R)^{-1}$; $\hat{\mathbf{x}}_{k} = \hat{\mathbf{x}}_{k|k-1} + K(z_k - H\hat{\mathbf{x}}_{k|k-1})$
    \end{itemize}
    $F$: state transition; $H$: observation model; $Q$, $R$: process/measurement noise covariances.
\visualseparator
    \begin{center}
    \begin{tikzpicture}[
        box/.style={rectangle, draw, rounded corners=2pt, minimum width=2.2cm, minimum height=0.7cm,
                    font=\scriptsize, fill=blue!10, align=center},
        arr/.style={-{Stealth[length=1.5mm]}, thick}
    ]
        \node[box] (pred) {Predict\\$\hat{x}_{k|k-1}, P_{k|k-1}$};
        \node[box, right=1.8cm of pred] (upd) {Update\\$\hat{x}_k, P_k$};
        \node[box, below=0.8cm of upd, fill=orange!15] (meas) {Measurement\\$z_k$};
        \draw[arr] (pred) -- node[above, font=\tiny] {prior} (upd);
        \draw[arr] (meas) -- (upd);
        \draw[arr] (upd.east) -- ++(0.6,0) |- node[right, font=\tiny, pos=0.25] {next step} (pred.east);
    \end{tikzpicture}
\captionof{figure}{The Kalman filter's two alternating steps. Prediction moves
the estimate forward and \emph{increases} uncertainty; the update folds in a
measurement and decreases it. The filter is optimal only while its linear and
Gaussian assumptions hold.}
\label{fig:cv-kalman-cycle}
    \end{center}

\end{v3topic}
\begin{v3topic}{Particle Filters and Deep Tracking}
    \textbf{Particle filter} (Sequential Monte Carlo): represents the posterior by $N$ weighted samples (particles)\textsuperscript{\cite[][pp.~548--553]{szeliski2022}}.
    Steps: (1) propagate via motion model, (2) weight by observation likelihood, (3) resample.
    Handles non-linear, non-Gaussian distributions; computationally heavy for large $N$.
\visualseparator
    \textbf{Deep tracking}:
    \begin{itemize}
        \item \textbf{SORT} (2016): Kalman filter state $(x,y,w,h,\dot{x},\dot{y})$ + Hungarian matching via IoU distance; fast, simple.
        \item \textbf{DeepSORT} (2017): adds deep ReID appearance descriptor to reduce ID switches.
        \item \textbf{ByteTrack} (2022): retains low-confidence detections in association; state-of-the-art MOT performance.
    \end{itemize}
\end{v3topic}

% ---- 9.2.5 Shape ------------------------------------------------------------------------------
\subsection{Shape}

\begin{v3topic}{Shape from Shading and Stereo}
    \textbf{Shape from Shading}: recover surface normals from a single image given known illumination\textsuperscript{\cite[][pp.~667--676]{szeliski2022}}.
    Under the Lambertian model: $I = \rho\,(\mathbf{n} \cdot \mathbf{l})$
    where $\rho$ is albedo, $\mathbf{n}$ surface normal, $\mathbf{l}$ light direction.
    The shape-albedo ambiguity makes this under-constrained without additional priors.
    \textbf{Photometric stereo}: multiple images under different illumination directions; resolves ambiguity.
\visualseparator
    \textbf{Active contours / snakes} (Kass et al., 1988): energy-minimising curve
    $E = \int \left[\alpha|C'|^2 + \beta|C''|^2 + E_{\mathrm{image}}(C)\right] \mathrm{d}s$
    (internal smoothness + external image gradient). Level-set methods provide an implicit alternative that handles topology changes.

\end{v3topic}
\begin{v3topic}{3D Reconstruction and Structure from Motion}
    \textbf{Structure from Motion (SfM)} recovers a sparse 3D point cloud and camera poses from an unordered image set\textsuperscript{\cite[][pp.~719--733]{szeliski2022}}:
    \begin{enumerate}
        \item Feature extraction and matching (SIFT + RANSAC)
        \item Incremental or global camera pose estimation
        \item Triangulation of 3D points
        \item \textbf{Bundle adjustment}: non-linear least-squares refining all camera parameters and point positions simultaneously
    \end{enumerate}
    \textbf{Multi-View Stereo (MVS)}: extends SfM to dense reconstruction.
    Modern neural approaches (NeRF, 3DGS) learn implicit or explicit 3D representations directly from images via differentiable rendering.
\end{v3topic}


% ==== Section 3: Computer Vision for Autonomous Systems =====================================
\section{Computer Vision for Autonomous Systems}

% ---- 9.3.1 Image Formation and Acquisition -------------------------------------------------
\subsection{Image Formation and Acquisition}

\begin{v3topic}{Light and Color Models}
    Cameras sample a 3D scene illuminated by light (visible spectrum: 380--700\,nm) and project it onto a 2D sensor\textsuperscript{\cite[][pp.~15--20]{jahneDigitalImageProcessing2005}}.
    Radiometry describes light transport (irradiance, radiance, BRDF); photometry weights by human perception.
\visualseparator
    \begin{tblr}{width=\linewidth,colspec={|Q[l,m]|X[l]|}}
        \hline
        \textbf{Color Model} & \textbf{Use Case} \\
        \hline
        RGB & Display, rendering (device-dependent) \\
        \hline
        HSV/HSL & Color-based segmentation, skin detection \\
        \hline
        CIE LAB & Perceptually uniform; color comparison \\
        \hline
        YCbCr & Video compression (JPEG, H.264) \\
        \hline
    \end{tblr}
\captionof{table}{Colour models by use case: additive for displays, subtractive for print, perceptual for analysis.}
\label{tab:cv-colour-models}

\end{v3topic}
\begin{v3topic}{Camera Calibration}
    \textbf{Zhang's method} (2000) estimates camera intrinsics $K$, extrinsics $[R|t]$, and lens distortion from multiple views of a planar checkerboard\textsuperscript{\cite[][pp.~45--52]{howseLearningOpenCV42020}}.
    Each view provides 2 constraints on $K$; $\geq 3$ views are needed.
    Radial distortion is modelled as $r'(x) = x(1 + k_1 r^2 + k_2 r^4 + \ldots)$.
    Implemented in OpenCV via \texttt{cv2.calibrateCamera()}.
\end{v3topic}

% ---- 9.3.2 Sensors -------------------------------------------------------------------------
\subsection{Sensors for Autonomous Systems}

\begin{v3topic}{Vision Sensors}
    \begin{itemize}
        \item \textbf{Monocular camera}: single image; depth ambiguous; cheap, widely deployed
        \item \textbf{Stereo cameras}: depth from disparity $d = fB/Z$ (baseline $B$, focal length $f$)
        \item \textbf{Event cameras} (DVS): log changes in illuminance asynchronously per pixel; microsecond latency, high dynamic range, no motion blur
        \item \textbf{Night vision}: near-infrared (NIR) illumination + sensor, or thermal (LWIR, 8--14\,$\mu$m)
    \end{itemize}

\end{v3topic}
\begin{v3topic}{LiDAR and Radar}
    \textbf{LiDAR} measures round-trip time of laser pulses: $d = ct/2$\textsuperscript{\cite[][pp.~8--10]{howseLearningOpenCV42020}}.
    \begin{itemize}
        \item \emph{Rotating} (Velodyne): 360\textdegree{} scan; 16--128 beams; dense point cloud at 10\,Hz
        \item \emph{Solid-state} (MEMS/OPA/Flash): no moving parts; cheaper; narrower FoV
    \end{itemize}
    \textbf{Radar} (FMCW, 77\,GHz): measures range \emph{and} velocity simultaneously via Doppler. All-weather capable; sparse but robust.

\end{v3topic}
\begin{v3topic}{Sensor Fusion}
    Autonomous systems combine complementary sensors to overcome individual limitations:
    cameras provide rich texture/colour, LiDAR provides accurate 3D geometry, radar provides velocity and weather robustness.
\visualseparator
    \begin{tblr}{width=\linewidth,colspec={|Q[l,m]|X[l]|}}
        \hline
        \textbf{Fusion Level} & \textbf{Description} \\
        \hline
        Early (raw) & Combine raw sensor data (e.g., project LiDAR into image) \\
        \hline
        Middle (feature) & Fuse intermediate feature representations from each modality \\
        \hline
        Late (decision) & Combine per-sensor detections/decisions \\
        \hline
    \end{tblr}
\captionof{table}{Fusion levels. Fusing earlier preserves more information and costs more bandwidth and calibration effort.}
\label{tab:cv-sensor-fusion}
    State estimation typically uses \textbf{Kalman filters} (linear) or \textbf{particle filters} (non-linear, non-Gaussian) to fuse temporal sensor streams.
\end{v3topic}

% ---- 9.3.3 Object Detection and Tracking ---------------------------------------------------
\subsection{Object Detection and Tracking}

\begin{v3topic}{Classical Object Detection}
    Before deep learning, detection relied on hand-crafted features + classifiers:
    \begin{itemize}
        \item \textbf{Viola-Jones} (2001): Haar-like features + AdaBoost cascade; first real-time face detector
        \item \textbf{HOG + SVM} (Dalal \& Triggs, 2005): histogram of oriented gradients in cells + linear SVM; pedestrian detection standard
        \item \textbf{Sliding window}: evaluate classifier at every position/scale; exhaustive but slow
    \end{itemize}

\end{v3topic}
\begin{workedexample}[Intersection over union, and why the threshold matters]
Boxes are $(x_1,y_1,x_2,y_2)$. Ground truth $B = (30,30,70,70)$, area
$40\times40 = 1600$.

\emph{Case 1 --- a loose prediction} $A = (10,10,50,50)$, area $1600$.
\[
\begin{aligned}
  \text{overlap in }x &: \max(10,30)=30 \text{ to } \min(50,70)=50
    \Rightarrow 20,\\
  \text{overlap in }y &: 30 \text{ to } 50 \Rightarrow 20,\\
  \text{intersection} &= 20\times20 = 400,\\
  \text{union} &= 1600 + 1600 - 400 = 2800,\\
  \mathrm{IoU} &= 400/2800 = \mathbf{0.14}.
\end{aligned}
\]
Below the usual $0.5$ threshold, so this counts as a \emph{false positive} ---
and the object also counts as \emph{missed}. One sloppy box costs twice.

\emph{Case 2 --- a tight prediction} $A = (15,15,55,55)$:
\[
  \text{intersection} = 35\times35 = 1225,\qquad
  \text{union} = 1600+1600-1225 = 1975,\qquad
  \mathrm{IoU} = \mathbf{0.62}.
\]
This passes at $\mathrm{IoU}\ge0.5$ and \emph{fails} at $\mathrm{IoU}\ge0.75$.

\emph{The lesson.} The same detector, unchanged, scores very differently
depending on a threshold chosen by the benchmark. That is why modern reporting
uses mAP averaged over thresholds from $0.50$ to $0.95$ --- and why a paper
quoting mAP@0.5 alone is quoting the most flattering number available.
\end{workedexample}

\begin{figure}[htbp]
\centering
\begin{tikzpicture}[x=0.55mm,y=0.55mm,
  gt/.style={draw=AlmanachTeal,very thick,fill=AlmanachTeal!12},
  pr/.style={draw=AlmanachOchre,very thick,dashed,fill=AlmanachOchre!12},
  lbl/.style={font=\sffamily\scriptsize,color=AlmanachInk!75}]
  % --- case 1: loose prediction ---
  \draw[gt] (30,30) rectangle (70,70);
  \draw[pr] (10,10) rectangle (50,50);
  \fill[AlmanachRed!35] (30,30) rectangle (50,50);
  \node[lbl] at (40,-14) {$\mathrm{IoU}=0.14$ --- rejected};
  % --- case 2: tight prediction ---
  \begin{scope}[xshift=55mm]
    \draw[gt] (30,30) rectangle (70,70);
    \draw[pr] (15,15) rectangle (55,55);
    \fill[AlmanachRed!35] (30,30) rectangle (55,55);
    \node[lbl] at (40,-14) {$\mathrm{IoU}=0.62$ --- accepted at 0.5};
  \end{scope}
  % --- legend ---
  \node[lbl,color=AlmanachTeal,anchor=west]  at (215,62) {ground truth};
  \node[lbl,color=AlmanachOchre,anchor=west] at (215,47) {prediction};
  \node[lbl,color=AlmanachRed,anchor=west]   at (215,32) {intersection};
\end{tikzpicture}
\caption{Two predictions against the same ground-truth box. Intersection over
union is the ratio of the shaded overlap to the total area covered by both; a
shift of only $20\,\%$ of the box width moves the score across the standard
acceptance threshold.}
\label{fig:cv-iou}
\end{figure}

\begin{v3topic}{Deep Object Detection}
    Deep detectors learn features and localisation end-to-end\textsuperscript{\cite[][pp.~105--109]{fleuretLittleBookDeep}}:
\visualseparator
    \begin{tblr}{width=\linewidth,colspec={|Q[l,m]|Q[c,m]|X[l]|}}
        \hline
        \textbf{Model} & \textbf{Type} & \textbf{Key Innovation} \\
        \hline
        Faster R-CNN & Two-stage & Region Proposal Network + ROI pooling \\
        \hline
        SSD & Single-stage & Multi-scale feature maps; anchor boxes \\
        \hline
        YOLO & Single-stage & Single forward pass; grid-based; real-time \\
        \hline
        DETR & Transformer & Bipartite matching; no anchors or NMS \\
        \hline
    \end{tblr}
\captionof{table}{Deep detectors by family. Two-stage detectors trade latency for accuracy; single-stage detectors made real-time detection practical.}
\label{tab:cv-detectors}
    Key concepts: \textbf{IoU} (Intersection over Union), \textbf{NMS} (Non-Maximum Suppression), \textbf{mAP} (mean Average Precision).

\end{v3topic}
\begin{v3topic}{Multi-Object Tracking}
    Tracking extends detection across time by associating detections frame-to-frame:
    \begin{itemize}
        \item \textbf{SORT} (2016): Kalman filter prediction + Hungarian algorithm for bounding box association; uses IoU as distance metric
        \item \textbf{DeepSORT} (2017): adds a deep ReID appearance descriptor; reduces ID switches by re-identifying occluded targets
        \item \textbf{ByteTrack} (2022): uses \emph{low-confidence} detections for association; SOTA on MOT17
    \end{itemize}
    Evaluation metrics: \textbf{MOTA} (tracking accuracy), \textbf{MOTP} (precision), \textbf{ID switches} (identity changes).
\end{v3topic}

% ---- 9.3.4 Semantic Segmentation -----------------------------------------------------------
\subsection{Semantic and Instance Segmentation}

\begin{v3topic}{Semantic Segmentation}
    Assigns a class label to \emph{every pixel} in the image\textsuperscript{\cite[][pp.~110--112]{fleuretLittleBookDeep}}.
    \begin{itemize}
        \item \textbf{FCN} (Long et al., 2015): fully convolutional; transposed convolutions for upsampling; first end-to-end pixel labelling
        \item \textbf{U-Net} (Ronneberger et al., 2015): encoder--decoder with skip connections; standard for medical imaging
        \item \textbf{DeepLab v3+} (Chen et al., 2018): dilated (atrous) convolutions + ASPP for multi-scale context
    \end{itemize}

\end{v3topic}
\begin{v3topic}{Instance and Panoptic Segmentation}
    \textbf{Instance segmentation} distinguishes individual objects of the same class:
    \begin{itemize}
        \item \textbf{Mask R-CNN} (He et al., 2017): Faster R-CNN + per-instance mask branch; ROI Align for pixel-accurate alignment
    \end{itemize}
    \textbf{Panoptic segmentation} unifies \emph{stuff} (amorphous regions: sky, road) and \emph{things} (countable objects: cars, people) --- every pixel gets a class label \emph{and} an instance ID.
    \textbf{MOTS} (Multi-Object Tracking and Segmentation) extends this to video with per-frame masks and tracking IDs.
\end{v3topic}


% ==== Section 4: Cognitive Computer Vision ===================================================
\section{Cognitive Computer Vision}

% ---- 9.4.1 Image Classification ------------------------------------------------------------
\subsection{Image Classification}

\begin{v3topic}{CNN Classifiers and Transfer Learning}
    Image classification assigns one or more class labels to an entire image\textsuperscript{\cite[][p.~104]{fleuretLittleBookDeep}}.
    Standard architectures (AlexNet, VGG, ResNet --- see Ch.~6) are trained on ImageNet (1.2M images, 1000 classes) with cross-entropy loss.
    \textbf{Transfer learning}: pre-train on ImageNet, freeze early layers, fine-tune later layers on the target dataset --- dramatically reduces labelled data requirements.

\end{v3topic}
\begin{v3topic}{Image Retrieval and Scene Understanding}
    \textbf{Image retrieval}: extract a CNN embedding (e.g., penultimate layer) and find nearest neighbours via cosine similarity in a vector index (FAISS)\textsuperscript{\cite[][p.~104]{fleuretLittleBookDeep}}.
    \textbf{Scene understanding} goes beyond object labels: recognising spatial layout, relationships between objects, and functional properties of a scene.
    Modern approaches use \textbf{Vision Transformers (ViT)}: split image into patches, treat as tokens, and apply self-attention --- competitive with CNNs on classification.
\end{v3topic}

% ---- 9.4.2 Recognition ---------------------------------------------------------------------
\subsection{Face Recognition}

\begin{v3topic}{Face Recognition Pipeline}
    Face recognition follows a standard pipeline: (1) \textbf{detection} (locate face in image), (2) \textbf{alignment} (geometric normalisation), (3) \textbf{feature extraction} (embed into vector), (4) \textbf{comparison} (distance in embedding space).
    \begin{itemize}
        \item \textbf{Eigenfaces} (Turk \& Pentland, 1991): PCA-based global appearance; limited by illumination/pose
        \item \textbf{FaceNet} (Schroff et al., 2015): triplet loss $\to$ 128-d embedding; L2 distance for verification
        \item \textbf{ArcFace} (Deng et al., 2019): additive angular margin loss; SOTA on LFW, MegaFace
    \end{itemize}

\end{v3topic}
\begin{v3topic}{Verification vs.\ Identification}
    \begin{tblr}{width=\linewidth,colspec={|Q[l,m]|X[l]|X[l]|}}
        \hline
        \textbf{Task} & \textbf{Question} & \textbf{Complexity} \\
        \hline
        Verification & ``Is this person $X$?'' (1:1) & Compare against one template \\
        \hline
        Identification & ``Who is this?'' (1:$N$) & Search entire gallery; harder \\
        \hline
    \end{tblr}
\captionof{table}{Verification asks a one-to-one question and identification a one-to-many one --- which is why identification error rates grow with gallery size.}
\label{tab:cv-verification-identification}
    Different threat models: verification tolerates higher thresholds (fewer false accepts); identification must handle large galleries efficiently.
\end{v3topic}

% ---- 9.4.3 Image Synthesis -----------------------------------------------------------------
\subsection{Image Synthesis}

\begin{v3topic}{Super-Resolution and Style Transfer}
    \textbf{Super-resolution}: reconstruct high-resolution images from low-resolution input.
    SRCNN (Dong et al., 2014) was the first deep approach; ESRGAN (Wang et al., 2018) uses GAN-based perceptual loss for photorealistic results.

    \textbf{Neural style transfer} (Gatys et al., 2015): optimise an image to match the \emph{content} of one image (via feature activations) and the \emph{style} of another (via Gram matrix of feature maps).

\end{v3topic}
\begin{v3topic}{Image-to-Image Translation}
    \begin{itemize}
        \item \textbf{pix2pix} (Isola et al., 2017): conditional GAN with \emph{paired} training data; learns mapping between image domains (sketch $\to$ photo, segmentation $\to$ image)
        \item \textbf{CycleGAN} (Zhu et al., 2017): \emph{unpaired} domain translation using cycle-consistency loss ($G(F(x)) \approx x$); enables horse $\leftrightarrow$ zebra, summer $\leftrightarrow$ winter
    \end{itemize}
    Both demonstrate that GANs (see Ch.~6) can learn complex image transformations beyond simple generation.
\end{v3topic}

% ---- 9.4.4 Vision and Language --------------------------------------------------------------
\subsection{Vision and Language}

\begin{v3topic}{CLIP and Zero-Shot Vision}
    \textbf{CLIP} (Radford et al., 2021) jointly trains an image encoder and a text encoder on 400M image--text pairs using contrastive learning\textsuperscript{\cite[][pp.~115--117]{fleuretLittleBookDeep}}.
    Similarity: $\ell_{m,n} = f(i_m)^\top g(t_n)$.
    This enables \textbf{zero-shot classification}: compare image embedding against text embeddings of class descriptions --- no task-specific training needed.

\end{v3topic}
\begin{v3topic}{Image Captioning and VQA}
    \textbf{Image captioning}: CNN image encoder $\to$ transformer/LSTM text decoder; generates natural language descriptions of images.
    \textbf{Visual Question Answering (VQA)}: given an image and a text question, predict an answer --- requires joint vision--language reasoning.

\end{v3topic}
\begin{v3topic}{Text-to-Image Generation}
    Text-conditioned image generation has progressed rapidly:
    \begin{itemize}
        \item \textbf{DALL-E} (Ramesh et al., 2021): discrete VAE + autoregressive transformer; generates images from text prompts
        \item \textbf{Stable Diffusion} (Rombach et al., 2022): \emph{latent diffusion} --- denoising in a compressed latent space; text conditioning via CLIP encoder; open-source, runs on consumer GPUs
    \end{itemize}
    The \textbf{diffusion process}\textsuperscript{\cite[][pp.~121--123]{fleuretLittleBookDeep}}: a forward process gradually adds Gaussian noise to data; the model learns to reverse this, generating samples by iterative denoising from pure noise.
\end{v3topic}


% ==== Section 5: Challenges in Computer Vision ===============================================
\section{Challenges in Computer Vision}

% ---- 9.5.1 Fairness and Explainability ------------------------------------------------------
\subsection{Fairness and Explainability}

\begin{v3topic}{Bias in Computer Vision}
    Face recognition error rates vary significantly by demographic group (the \emph{Gender Shades} study by Buolamwini \& Gebru, 2018, found error rates up to 34$\times$ higher for dark-skinned women than light-skinned men).
    Bias arises from unbalanced training data, label quality, and evaluation practices.
    Mitigations: balanced datasets, fairness-aware training objectives, demographic-specific evaluation.

\end{v3topic}
\begin{v3topic}{Explainability Methods}
    Understanding \emph{why} a model makes a decision is essential for trust and debugging:
    \begin{itemize}
        \item \textbf{Grad-CAM} (Selvaraju et al., 2017): gradient of class score w.r.t.\ final conv feature maps $\to$ localisation heatmap
        \item \textbf{Saliency maps}: backpropagate class score to input pixels to show which pixels matter
        \item \textbf{LIME}: model-agnostic; perturb input locally, fit a simple interpretable model
    \end{itemize}
\end{v3topic}

% ---- 9.5.2 Self-Supervised and Contrastive Learning -----------------------------------------
\subsection{Self-Supervised and Contrastive Learning}

\begin{v3topic}{Learning Without Labels}
    Self-supervised methods learn visual representations from unlabelled data, then transfer to downstream tasks:
\visualseparator
    \begin{tblr}{width=\linewidth,colspec={|Q[l,m]|Q[c,m]|X[l]|}}
        \hline
        \textbf{Method} & \textbf{Year} & \textbf{Key Idea} \\
        \hline
        SimCLR & 2020 & Two augmented views + NT-Xent contrastive loss; large batch needed \\
        \hline
        BYOL & 2020 & Online + target networks; stop-gradient; no negative pairs \\
        \hline
        DINO & 2021 & ViT + self-distillation; student on local, teacher on global crops \\
        \hline
        MAE & 2021 & Mask 75\% of image patches; reconstruct; efficient scalable pre-training \\
        \hline
    \end{tblr}
\captionof{table}{Self-supervised methods and the pretext task each uses to manufacture a training signal from unlabelled data.}
\label{tab:cv-self-supervised}
    These methods approach or surpass supervised pre-training on ImageNet, especially when labelled data is scarce.
\end{v3topic}

% ---- 9.5.3 Robust Computer Vision -----------------------------------------------------------
\subsection{Robust Computer Vision}

\begin{v3topic}{Adversarial Examples}
    Imperceptible perturbations can cause confident misclassification (Szegedy et al., 2013; Goodfellow et al., 2015).
    \textbf{FGSM} (Fast Gradient Sign Method):
    \begin{equation}
        x_{\mathrm{adv}} = x + \epsilon \cdot \mathrm{sign}\!\left(\nabla_x \mathcal{L}(\theta, x, y)\right)
    \end{equation}
\where{$x$ is the clean input, $y$ its true label, $\theta$ the model
parameters, and $\epsilon$ the perturbation budget --- deliberately small enough
that $x_{\mathrm{adv}}$ looks unchanged to a person.}
    Defences: adversarial training (augment with adversarial examples), certified robustness, input preprocessing.

\end{v3topic}
\begin{v3topic}{Domain Shift and OOD Detection}
    \textbf{Domain shift}: distribution mismatch between training and deployment environments (e.g., lab images vs.\ real-world scenes).
    Mitigation: domain adaptation via adversarial alignment of feature distributions, or test-time adaptation.

    \textbf{Out-of-Distribution (OOD) detection}: identify inputs that fall outside the training distribution.
    Methods: maximum softmax probability, energy-based scoring, Mahalanobis distance in feature space.
    Critical for safety in autonomous driving and medical imaging.
\end{v3topic}


% ==== Section 6: Further Reading =============================================================
%
\begin{workedexample}[Why three $3\times3$ layers beat one $7\times7$]
\emph{Receptive field.} Stacking convolutions with stride $1$ grows the
receptive field additively: each $k\times k$ layer adds $k-1$.
\[
  \text{after 1 layer: } 3,\qquad
  \text{after 2: } 3+2 = 5,\qquad
  \text{after 3: } 5+2 = 7 .
\]
Three $3\times3$ layers see exactly as much input as one $7\times7$ layer.

\emph{Parameters}, for $C$ input and $C$ output channels:
\[
  3\times(3\times3\times C\times C) = 27C^2
  \quad\text{versus}\quad
  7\times7\times C\times C = 49C^2 .
\]
That is $45\,\%$ fewer parameters --- and two extra nonlinearities, so the
stack is also strictly more expressive. This single comparison is VGG's entire
design argument\textsuperscript{\cite{simonyanVeryDeepConvolutional2015}}, and
it is why almost every convolutional architecture since uses small kernels.
\end{workedexample}

\begin{workedexample}[Activations, not weights, decide your batch size]
Train on batches of $32$ RGB images at $224\times224$.

\emph{The input tensor:}
\[
  32 \times 224 \times 224 \times 3 \times 4\,\text{bytes}
  = 19.3\,\text{MB}.
\]
Negligible.

\emph{One early activation map} with $64$ channels at full resolution:
\[
  32 \times 224\times224 \times 64 \times 4\,\text{bytes}
  = 411\,\text{MB} \quad\text{--- for a single layer.}
\]
And backpropagation must \emph{keep} these activations for the backward pass
(\cref{sec:dl-backprop}), across every layer.

\emph{The consequence.} A network with $25$ million parameters occupies about
$100\,\text{MB}$ of weights, while its stored activations at this batch size
run to many gigabytes. When you hit an out-of-memory error, the fix is almost
never a smaller model --- it is a smaller batch, a lower input resolution, or
activation checkpointing, which trades roughly $30\,\%$ extra compute for a
large memory saving.
\end{workedexample}

\begin{cautionbox}[Vision models latch onto the background]
A classifier trained on photographs frequently learns the context rather than
the object: cows on grass, boats on water, a ruler beside a skin lesion, a
particular hospital's scanner watermark. Accuracy on a held-out split from the
same source looks excellent, because the shortcut is present there too, and
performance collapses the moment the context changes.

Three cheap checks: evaluate on data from a genuinely different source; occlude
the object and confirm the prediction degrades; and inspect saliency or occlusion
maps on a sample of correct predictions, not just the errors. A model that is
right for the wrong reason will be right until it is catastrophically wrong.
\end{cautionbox}

\begin{practicallab}[Find out what your vision model is actually looking at]
\textbf{Goal.} Replace ``the model achieves 94\,\% accuracy'' with a statement
about \emph{why}.

\textbf{Protocol.}
\begin{enumerate}
  \item Take any pretrained classifier and a set of images it labels correctly
        with high confidence.
  \item \textbf{Occlusion test.} Slide a grey square across each image and
        record how the confidence for the true class changes. Plot the result
        as a heat map. Where does the evidence live?
  \item \textbf{Background swap.} Segment the object and paste it onto three
        unrelated backgrounds. Does the prediction survive?
  \item \textbf{Corruption sweep.} Apply Gaussian noise, blur, JPEG
        compression, and brightness shift at five severities each. Plot
        accuracy against severity.
  \item \textbf{Resolution sweep.} Downscale until accuracy collapses. Compare
        the pixel height at that point with the pinhole calculation earlier in
        this chapter.
\end{enumerate}

\textbf{Deliverable.} One heat map per class, one accuracy-versus-severity plot,
and a sentence naming the corruption your model is least robust to.
\textbf{The point:} step 3 regularly demotes an apparently excellent model, and
it takes an afternoon.
\end{practicallab}

\begin{checkpoint}
Your pedestrian detector scores mAP@0.5 of $0.91$ on the test set and fails in
deployment at dusk. Using this chapter, list the three measurements you would
take first --- one about the sensor, one about the evaluation protocol, and one
about distribution shift --- and say what each would rule in or out.
\end{checkpoint}

\section{Deep Dive: Foundation Vision and Robust Perception}
\begin{v3topic}{Prompts Change the Vision Interface}
Promptable segmentation and vision--language models replace a fixed class list
with points, boxes, text, or examples.
This flexibility expands evaluation: test localisation, open-vocabulary
recognition, grounding, refusal, and consistency across prompt wording.

\end{v3topic}
\begin{v3topic}{Shift Is Often Physical}
Weather, optics, compression, sensor ageing, viewpoint, motion blur, and scene
composition alter the image-generating process.
Corruption benchmarks are useful probes but do not replace data collected in
the intended operational domain.

\end{v3topic}
\begin{v3topic}{Perception-to-Action Trace}
For one detected object, trace calibration, preprocessing, model output,
tracking, fusion, uncertainty, planning, and fallback.
Express how spatial or temporal error propagates into an action margin.
Evaluate the system-level consequence, not only IoU or mAP in isolation.
\end{v3topic}

\chapterreview{Sense}{Represent}{Infer}{Stress-test}

\section{Further Reading}

\begin{v3topic}{Recommended Sources for Computer Vision}
    \begin{itemize}
        \item \fullcite{szeliski2022} --- Comprehensive CV textbook covering geometry, image processing, feature detection, recognition, and 3D reconstruction.
        \item \fullcite{gonzalezWoods2018} --- Classic digital image processing reference: filtering, morphology, frequency domain, segmentation.
        \item \fullcite{hartleyZisserman2004} --- Definitive treatment of multi-view geometry: epipolar geometry, fundamental matrix, 3D reconstruction.
        \item \fullcite{jahneDigitalImageProcessing2005} --- Thorough coverage of image processing fundamentals with mathematical rigour.
        \item \fullcite{howseLearningOpenCV42020} --- Practical computer vision with OpenCV and Python: camera calibration, feature detection, tracking.
        \item \fullcite{fleuretLittleBookDeep} --- Concise deep learning overview including CNNs for vision, object detection, and segmentation architectures.
    \end{itemize}
\end{v3topic}
